\section{Data and method}\label{sec:method}

Three of the design decisions below were not ours to make. Each was imposed by the
data-processing authorisation under which the study was conducted, an instrument issued by the
administration acting as data controller, and each is stated where it arises rather than
justified after the fact. The distinction matters and we hold to it throughout: the
\emph{obligation} in each case is contractual, arising from that instrument; the statutes
supply the standards the instrument implements and the interests it protects. We do not claim
that any statute independently compelled these choices.

One property of the design carries the auditability claim, and it is a property of the
artefacts rather than of this prose: every stage declares the question it answers, and that
declaration is copied into each artefact the stage writes, beside the code version, the
versions of the numerical libraries, the resolved parameters and a digest of every input
consumed. A reader who disputes a conclusion can locate the artefact that produced it, read the
question it claims to answer, and re-run the stage over the same inputs.
Figure~\ref{fig:pipeline} lays out the whole route; the subsections below follow it stage by
stage.

\subsection{Corpus}\label{sec:corpus}

The corpus is drawn from the electronic process system of a state-level secretariat of a
Brazilian state administration, read through the administration's analytical mirror of that
system rather than through its transactional database. The mirror holds 21{,}063{,}844 recorded
activities between January 2021 and the extraction date, so restricting the study to one
calendar year is a design decision and not a limit of the source.

The entry predicate is \emph{born in the body}, and the distinction is not cosmetic: a case
counts when its creation event was executed by a unit of that secretariat inside the year, not
when it merely passed through, and not by the filing field the source system records. A case
enters the analysis if it satisfies that predicate, concluded within the same year, and has a
retrievable opening document. The full funnel matters more than its last step and we give it
here: \textbf{12{,}621} cases were born at that secretariat during the year; \textbf{5{,}445}
of them also concluded within it and make up the loaded corpus; and \textbf{5{,}243} of those
carry a readable opening document and enter the analysis. The 7{,}176 cases that had not
concluded by 31 December are excluded by construction rather than by estimate, and they are
structurally the long ones. Of the documents that were sought and not obtained, 96 opening
documents were cancelled in the source system and can never be retrieved.

Two measures of ``concluded'' coexist in the sources and \textbf{cannot be added or presented
as one}. The graph's measure requires that no unit still hold the case open at the end of the
window; the mirror's is textual, recognising a conclusion from the description of the recorded
step. The second counts 7{,}920 conclusions logged inside the year against 4{,}701 cases that
logged none --- figures that answer a different question from the one the corpus answers, and
that are reported here so the difference stays visible rather than to be reconciled.

The corpus is pinned to what the graph actually loaded. Re-deriving it live from the curated
mirror returns 5{,}444 cases rather than 5{,}445, because a case concluded inside the year can
receive a later step and then fail the recency test the mirror applies. That re-derivation is a
drift check on a source under active curation, not a second measurement of the corpus, and
every downstream figure was computed on the pinned population.

What the corpus contains bears on how far the artefact reaches, and it should be stated plainly
rather than left to be inferred. The secretariat is responsible for personnel and general
administration, and its caseload is predominantly \emph{internal} administrative procedure ---
leave, intern placement, payroll deductions, contracting and the like. Some procedures cross
into other secretariats, other levels of government or municipalities. The party to a case is
therefore typically a public servant rather than a member of the public. Section~\ref{sec:legal}
states what follows for the normative claim; the short version is that the catalogue derived
here instantiates the service-charter duty only for that subset of procedures which are services
rendered to an external user, and the corpus is not principally composed of those.

Restricting to a single year was in part a cost decision for an initial study: five years of
event log would multiply ingestion, embedding, the clustering parameter sweeps and an
evaluation matrix that already runs for about eleven hours. The restriction is not only
economic, though. The authorisation delimits its object to cases born and moved at that
secretariat within the same twelve months, so the study window and the authorised window
coincide by construction.

\subsubsection{Right-censoring, measured}\label{sec:bias}

Requiring a case to open \emph{and} conclude within one calendar year truncates the sample in a
way that must be quantified rather than assumed benign. The time available to conclude shrinks
month by month: a case opened in January has 365 days, one opened in December has 31. The
observed maxima track that ceiling almost exactly --- 343 days against 365 available in January,
29 against 31 in December --- which is the signature of censoring rather than of how long these
procedures actually take.

%% GENERATED by `pipeline run results` - do not edit by hand.
\begin{table}[htbp]
\centering
\caption{Right-censoring by opening month.}\label{tab:corpus-bias}
\footnotesize
\begin{tabular}{lrrrrr}
\toprule
month & observed & days available & median life & max life & est.\ excluded \\
\midrule
2025-01 & 474 & 365 & 21 & 343 & 0.0\% \\
2025-02 & 619 & 334 & 14 & 290 & 0.2\% \\
2025-03 & 478 & 306 & 17 & 284 & 0.8\% \\
2025-04 & 470 & 275 & 25 & 244 & 2.1\% \\
2025-05 & 415 & 245 & 24 & 225 & 5.1\% \\
2025-06 & 621 & 214 & 18 & 204 & 21.7\% \\
2025-07 & 653 & 184 & 12 & 166 & 24.1\% \\
2025-08 & 398 & 153 & 18 & 129 & 26.2\% \\
2025-09 & 539 & 122 & 10 & 105 & 29.3\% \\
2025-10 & 381 & 92 & 10 & 75 & 31.9\% \\
2025-11 & 268 & 61 & 5 & 42 & 35.9\% \\
2025-12 & 129 & 31 & 5 & 29 & 44.5\% \\
\bottomrule
\end{tabular}
\par\vspace{2pt}\footnotesize\raggedright Observed 5,445; estimated true 6,658; estimated missing \textbf{1,213} (18.2\%). Lower bound: the reference month is itself truncated.
\end{table}


Taking the least-truncated month as a reference distribution, the number of cases that would
have concluded within a full-length window is estimated at 6{,}658 against the 5{,}445 observed,
so an estimated 1{,}213 cases --- 18.2\% of the \emph{concluded} population --- are missing, and
those are systematically the long ones. The scope of that figure should be stated precisely,
because it is easy to read as more than it is: it is a shortfall within the cases that conclude,
estimated from month-by-month censoring. It says nothing about the 7{,}176 cases that never
concluded inside the window at all, which are excluded outright. It is also a lower bound within
its own scope, since the reference month is itself truncated. Three consequences bind the rest of
the paper, and are restated in Section~\ref{sec:threats}: the corpus is not ``cases opened in one
year'' but ``cases opened in that year that also concluded quickly enough''; no duration or
efficiency statistic derived from it is admissible; and long, multi-unit procedures are
under-represented.

\subsection{Event knowledge graph}\label{sec:ekg}

Event data are modelled following established process-mining practice
\citep{vanderaalst2016processmining}. Extraction reads three sets from the analytical mirror ---
the recorded activities, the register of organisational units, and the metadata of each case ---
and applies the corpus predicate at read time, so nothing outside it is materialised.

The load proceeds in five chained steps, and the order between them is not indifferent.
Dimension nodes come first --- organisational units, users and declared procedure types --- so
that later loads merge into stable nodes instead of duplicating them. Activities are then
attached, each linked to its case, to the unit that executed it, to the user who performed it
and to its action type; case nodes are created from the case's own creation event, which fixes
one root per case before any activity hangs off it, and documents produced or referenced are
linked to the case that contains them. The third step is what turns a log into a process: the
activities of one case are chained by a directly-follows relation in recorded chronological
order \citep{esser2021ekg}. Holding quantities are derived over that chain --- how long each
unit retained each case and how often it received the case back, which is where rework becomes
visible --- and the status of each case at the end of the window is decided last, by the absence
of an open link in any unit rather than by a logged closing event.

The loaded graph holds 5{,}445 cases, 283{,}841 activities over 48 action types and 708
organisational units, and 42{,}658 documents that the event log references. Only one document
per case is read: its opening document. Deriving traces from a graph rather than a flat log
keeps the unit and action perspectives available simultaneously, which all three trace alphabets
of Section~\ref{sec:alphabets} draw on, and the same traversal yields the structural feature
blocks of Section~\ref{sec:features}. A dedicated stage then freezes that trace substrate to
disk, so the evaluation matrix and every later re-analysis run over exactly the same traces with
no database. We nevertheless chose a case notion --- the administrative case --- so the
divergence and convergence concerns that motivate object-centric process mining still apply
\citep{vanderaalst2019ocpm}.

\subsection{The opening document: reading, anonymisation, distillation}\label{sec:anon}

The opening document is converted to plain text first --- direct extraction where the file is
structured, reading by a locally executed vision model where it is not --- and suppression acts
on that text.

Personal data are pseudonymised or anonymised before any clustering or vectorisation step. This
ordering is not a robustness check: the governing instrument requires anonymisation to precede
those stages, and it invokes art.~12 of the Brazilian data protection law
\citep{brasil2018lgpd} in doing so. The reading matters. Article~12 is a
\emph{scope-delimiting} provision --- anonymised data fall outside the statute unless the
process can be reversed by reasonable means --- rather than a duty to anonymise. The duty here
is contractual; what art.~12 supplies is the standard against which the result is judged, and
the consequence that failing it returns the data to the statute's full scope. Precisely because
the transformation is imposed rather than chosen, whether it costs analytical utility is a
question the study must answer rather than assume, and Section~\ref{sec:res-rq5} answers it. The
entire pipeline is run twice, once on raw text and once on the anonymised variant, and the two
arms are compared case by case.

Detection combines two layers. The first is rule-based, in the spirit of engines built from
declarative recognisers \citep{microsoftpresidio}: national taxpayer identifiers for natural and
legal persons are recognised in punctuated and bare forms and \emph{validated by their own check
digits}, which stops protocol and document numbers from being suppressed by format resemblance;
e-mail addresses follow the standard character classes; and a staff registration number is
recognised only in context, preceded by the term that names it and above a minimum length. The
second layer is an open-schema neural entity extractor \citep{zaratiana2025gliner2}, whose spans
are mapped onto the same typed labels plus a catch-all for the identity and finance documents
the extractor emits under free-form keys --- an unmapped key would mean personal data left in
place, which is why the mapping is deliberately wide.

Replacement preserves referential integrity: the placeholder is derived from a digest of the
label and the entity in canonical form, so the same person receives the same placeholder in
every document in which they appear while the surface form disappears. Occurrences are replaced
longest span first, with each detector's own boundary guards --- digit lookarounds for the
numeric labels, word characters for the nominal ones --- because a single uniform guard would
silently miss the occurrences adjacent to digits. A verification mode re-runs the rule layer
\emph{over the output} and fails the run if any hit remains.

Removing names does not by itself anonymise an event log: a trace is itself a behavioural
fingerprint \citep{mannhardt2019privacy, elkoumy2022privacy}, and a sequence of dated moves
between named organisational units remains a behavioural record of the officials who handled
the case. We therefore state what the transformation does and does not cover, rather than
claiming anonymity as achieved.

\textbf{What it covers.} Direct and indirect identifiers in the free text of the opening
document --- personal names, national identifiers, addresses, contact details, registration
numbers --- are detected and replaced with typed placeholders before any embedding is computed.
\textbf{What it does not cover.} The trace itself is untouched: the sequence of organisational
units, the action types and their timestamps are all preserved, because they are the object of
study. A reader with access to the source system could re-identify a case from its routing
sequence alone. The arms therefore differ in the \emph{textual} channel only, which is the
scope within which the comparison in Section~\ref{sec:res-rq5} should be read. Calling the
result ``anonymised'' throughout is shorthand for this specific transformation, not a claim
that the statutory standard for anonymity has been met for the corpus as a whole.

From the treated text the method distils two shorter representations. An open-weight language
model produces, for each opening document, a \textbf{topic} --- the subject in a few words ---
and a one-sentence \textbf{summary}, both kept beside the text. The distillation exists because
the vector of the full text carries, alongside the subject, the administrative form the subject
arrives wrapped in; topic and summary are a subject signal with less of that noise, and they
enter the representation of Section~\ref{sec:features} as a block of their own. Because the
topic is derived from the text it differs between arms, so it is recomputed at every arm switch
rather than reused: reusing it would let the anonymised arm inherit the subject extracted from
raw text, and the comparison between arms would stop measuring what it claims to measure.

\subsection{Controlled keyword vocabulary}\label{sec:keywords}

From the same text comes an object of a different nature, and the only one in the whole method
a person can read: a controlled vocabulary of keywords. It is the bridge between the vector that
groups and the prose that names, and without it the naming step of Section~\ref{sec:naming}
regresses to document samples. The question it answers is which terms the collection itself uses
to describe its contents, without the list being written by hand, so the vocabulary is
\emph{growing} rather than fixed in advance: the language model proposes terms document by
document, and each proposed term either merges into a term already present or enters as new.

A term passes three chained filters. \textbf{Normalisation} lower-cases, strips accents and
punctuation, collapses whitespace, reduces unit paths to their final acronym and rejects merely
temporal expressions such as a bare year. \textbf{Two discard sets} follow --- the stopwords of
the language, and the terms that describe the administration itself: the body's own name, its
units, the vocabulary of case movement. Those appear in almost every document and therefore
separate nothing; keeping them would give the block a constant axis. \textbf{Variant merging}
comes last: two terms whose Dice similarity over character bigrams reaches 0.9 are treated as
one, and the first observed becomes the canonical form. A second pass re-annotates the documents
processed while the vocabulary was still being built, so no document is annotated against a
partial vocabulary. Over the 5{,}243 analysed cases, 2{,}297 controlled terms occur at least
once and 567 occur in at least three cases, which is the requirement for entering the feature
block and therefore fixes its width.

Three properties of that procedure are limitations rather than features. The similarity is
\textbf{lexical}, and errs in two predictable directions: it merges morphological variants of one
term, which is intended, but does not merge synonyms spelled differently and can merge short
terms that merely resemble each other. The choice of canonical form \textbf{depends on the
arrival order} of the documents, which makes the step non-deterministic across runs; the
vocabulary is therefore persisted as an artefact and reused, and it is the artefact rather than
the procedure that makes the dependent blocks reproducible. And the step must run \textbf{once,
before} both arms are built: run afterwards, the keyword block reaches vectorisation empty,
clustering proceeds over a null block and no check registers the absence.

\subsection{Features}\label{sec:features}

The substantive claim of the method sits in this step: procedural form and documentary content
are both signal, and neither alone identifies a process type. Each case is therefore represented
by six blocks that fuse what the case \emph{says} with how it \emph{moved}:

\begin{itemize}
  \item the opening document embedded with a multilingual sentence encoder \citep{chen2024bgem3} (1{,}024 dimensions);
  \item an LLM-produced topic and summary of that document, embedded the same way (1{,}024);
  \item binary TF--IDF over the organisational units a case visits (at most 500) and TF--IDF over
    consecutive unit pairs after collapsing repetitions (at most 1{,}000);
  \item a normalised histogram over the 48 action types;
  \item binary TF--IDF over the controlled keyword vocabulary of the opening document (567
    columns).
\end{itemize}

The two content blocks use \textbf{the opening document only}, and the restriction is
deliberate: later documents are mostly transmittal boilerplate that looks alike across thousands
of unrelated cases and dilutes the centroid, whereas the opening document is what the case
declares itself to be before the procedure has done anything to it. A case with no readable
document receives zero content blocks and is still grouped by its structural signature. Among the
structural blocks, the first records the presence of each unit on the route rather than how often
it appears; the second takes consecutive unit pairs after collapsing repetitions, so it describes
transition rather than dwell; and both are fitted over the batch itself, so a term's weight
reflects the set being clustered.

Each block is $L^2$-normalised per row before concatenation and all six carry weight 1.0, so the
scheme weights blocks equally regardless of their width: the 48-bin action histogram contributes
as much mass as the 1{,}024-dimensional document embedding. Every weight is exposed as a
parameter, and that is what makes the fusion claim testable by the ablation of
Section~\ref{sec:res-ablation}: a block can be zeroed with nothing else changing.

The concatenated vector, of several thousand dimensions, is reduced to \textbf{50} by UMAP
\citep{mcinnes2018umap} under cosine distance --- the metric that matches a space of normalised
embeddings --- with 15 neighbours and a fixed seed. The reduction reproduces bit for bit on the
same machine but not across processors and library versions; the stage declares that condition in
the artefact it writes, beside the seed, so a reader knows in advance which quantities reproduce
exactly and which reproduce only in distribution.

Both the language model that produces topics and keywords and the embedding model run
\emph{locally}, from open weights, on the administration's own high-performance computing
environment. This is a legal requirement rather than an engineering preference: the
authorisation confines processing to the administration's computing environment and forbids
external interfaces that would move case data outside it. A hosted commercial model was not an
option available to this study.

\subsection{Clustering}\label{sec:clustering}

\subsubsection{Three families, and why three}\label{sec:families}

Clustering runs with three algorithm families: density-based
\citep{campello2013hdbscan, mcinnes2017hdbscan}, centroid (\textit{k}-means) and modularity
\citep{traag2019leiden}. Running three is not methodological indecision: it is what separates a
property of the data from the inductive bias of one algorithm. Structure that appears under only
one of the three is a finding about that algorithm; structure that appears under all three is a
finding about the collection.

The three operate on the same 50-dimensional matrix, and parameters are swept rather than fixed.
The density family needs no declared number of clusters: it builds the density hierarchy of the
points and extracts the most stable partition from it, with a minimum cluster size of 20 and
excess-of-mass selection; a coarser variant merges clusters separated by less than
$\varepsilon = 0.5$. One implementation nuance binds the reading: the cosine metric describes the
\emph{reduction}, while density clustering runs in Euclidean distance over the already-reduced
space --- usual practice, since the projection preserves neighbourhood rather than angle, but not
to be read as cosine clustering. The modularity family operates over a $k$-nearest-neighbour
graph ($k = 15$, cosine) with the resolution swept and selected by cosine silhouette. The
centroid family sweeps $k$ over a grid with ten restarts per value and a fixed seed, under the
two stopping criteria described next. The sweep yields \textbf{seven} parameterisations per
textual arm --- the three families in a reference configuration plus four parameter variants ---
and therefore \textbf{fourteen} real clusterings across the two arms.

\subsubsection{Choosing \texorpdfstring{$k$}{k} and the resolution, swept and reported}\label{sec:selectk}

The number of clusters is not chosen by inspection. A dedicated stage sweeps the grid of $k$ and
reports, side by side, the elbow criterion --- defined as the maximum perpendicular distance to
the chord joining the endpoints of the inertia curve, with both axes normalised to the unit
interval, rather than read off a plot --- and the silhouette. Both values are carried forward as
distinct parameterisations instead of one being elected \emph{a priori}: the two criteria
optimise different objectives, compactness relative to separation in one case and marginal
inertia gain in the other, and before measuring there is no ground for preferring either. Carrying
both turns the choice of $k$ from a premise into an object of evaluation. Where computing the
silhouette over the full matrix exceeds a fixed time budget it switches to a fixed-size,
fixed-seed subsample, and the scope used is recorded in the artefact itself, so a subsample value
is never read as a full-matrix one. The whole sweep is persisted as an auditable artefact, not
only its winner.

\subsubsection{Density noise is assigned}\label{sec:noise}

The density family is the only one of the three that can refuse to assign a case to any cluster.
The refusal is informative, but it has a consequence that would make cross-family comparison
impossible: validity indices and downstream accuracy would be computed over \emph{different
populations}, and a family that evaluates only the easy cases is rewarded for it.

By default, therefore, each refused point is assigned to the cluster whose centroid is nearest,
in the same space in which the metrics are computed. Every family is then evaluated over the same
cases and the comparison between them becomes legitimate. The choice has a cost, which we state
rather than omit: a case that belongs to no family is filed under the nearest one, which is the
right trade for evaluation and the wrong one for a catalogue an official browses. What density
\emph{would} have refused is not lost --- it is recorded per clustering as the fraction refused
before assignment, and it is that quantity, rather than the final noise rate, that answers for
abstention in any analysis. Assigning to the nearest centroid is \emph{one} defensible choice
among others; the alternative of reclassifying the refused points into subclusters of their own
remains implemented and recoverable by parameter.

\subsubsection{One shuffled control per clustering}\label{sec:controls}

Each of the fourteen real clusterings is paired with a \textbf{control partition}: labels are
permuted across cases, preserving the cluster count, the size distribution and the evaluated
population exactly, and destroying only membership. It is against that control, and not against
a global model alone, that the gain is measured (Section~\ref{sec:estimand}).

The control is built \emph{per clustering} rather than as one shared null, for a substantive
reason: the shape of a partition --- how many clusters, of what sizes --- moves accuracy on its
own, independently of membership meaning anything. A control derived from a \emph{different}
clustering would measure a difference in granularity instead of a difference in grouping. The
control partition is therefore an artefact of the clustering phase, not of the evaluation phase.

\subsection{Naming}\label{sec:naming}

A catalogue entry needs a name, and the name is the part a reader sees first. Ours are generated
rather than authored, and naming runs in \textbf{two passes} --- a design that answers a failure
mode of the naive procedure. Naming each cluster in isolation from its most \emph{frequent}
keywords guarantees similar names for similar clusters, because within a family of neighbouring
procedures the most frequent terms are precisely the \emph{shared} ones, and a model that never
sees the competitors has no way to tell them apart.

In the first pass, therefore, each cluster reaches the locally executed language model with
\emph{two} keyword sets rather than one: the most frequent, which give the theme, and the
\textbf{distinctive} ones --- those over-represented in that cluster relative to the rest of the
collection, ranked by the ratio between the term's probability inside and outside the cluster,
with smoothing for rare terms. The latter separate; the former unite. Alongside them go a sample
of the most central cases --- the topic and summary of each opening document --- the
organisational units visited and the action types; and the model returns a name, a one-sentence
factual description and three distinguishing markers. It is instructed to anchor the name in the
cluster's own content and flow, to prefer a concrete subject over a generic label, and \emph{not}
to use the declared classification, so that the emergent type is not silently relabelled with the
taxonomy the study declines to rely on --- using it would make agreement against the declared
type circular, and that agreement is the external instrument by which the partition is assessed.

The second pass renames \emph{only} the clusters whose names collided, now with the rivals in
view: each receives the name, size and distinctive terms of its homonyms and an explicit
instruction to produce a name that separates them. Names whose embeddings remain near-duplicates
are merged by agglomerative clustering over name-embedding cosine similarity. Naming runs for
\emph{every} partition rather than only the one presumed to win, since the selection criterion is
applied afterwards and a partition without names would render an empty catalogue.

Two consequences follow and both are limitations rather than features. The names are unvalidated
machine output over administrative subject matter, which is a distinct artefact from the
partition itself and carries distinct risk; Section~\ref{sec:threats} states this alongside the
absence of specialist assessment of the types. And near-duplicate naming was an observed failure
mode during development rather than a hypothetical one, which is why the merge step exists.

\subsection{Trace alphabets}\label{sec:alphabets}

A case's trace is rendered in \textbf{three} alphabets, and the distinction delimits what each
number downstream can assert. The \textbf{action-by-unit} alphabet pairs each recorded action
with the unit that performed it --- the unit's full acronym, with no roll-up to a directorate ---
and is the raw state; the other two are projections of it, obtained by taking one of the two axes
and collapsing consecutive repetitions. The \textbf{routing} alphabet keeps only the sequence of
organisational units a case visits, abstracting what happened inside each; analysing the sequence
of units is an established organisational perspective \citep{song2008organizational}. The
\textbf{procedure} alphabet does the inverse, keeping the sequence of actions and abstracting the
organisational chart, which makes it the most transferable of the three because it does not
depend on this administration's unit names --- the generated tables label it the \emph{action}
alphabet, after the axis it retains.

Collapsing repetitions is the point of the two projections rather than an incidental loss: a case
that stays four events in one unit made \emph{one} visit, not four, and a unit that reappears
non-consecutively is a genuine return. Dwell does not leave the method --- the event count per
visit is kept beside the trace --- it only leaves the alphabet. Each trace also receives a
terminal sentinel state, without which a generator would have no way to learn \emph{when} a case
ends and would stop at the first conclusion state it met.

\subsubsection*{The alphabets are not compared with each other}

Each alphabet receives its own complete evaluation matrix, and the best model is elected
\emph{within} each. No number crosses alphabets at any point in this paper. The decision is one
of scope, and rests on three independent reasons:

\begin{enumerate}
  \item \textbf{The populations differ.} The two projections collapse consecutive repetitions, so
    a case handled entirely inside one unit reduces to a single routing state and is dropped for
    having no transition to predict. The routing alphabet therefore evaluates multi-unit cases,
    not the corpus.
  \item \textbf{The lengths differ.} Every measure sensitive to trace length --- and the suffix
    measures are --- changes scale between alphabets with nothing having changed in the
    collection.
  \item \textbf{The vocabularies differ by orders of magnitude.} A higher next-step accuracy over
    a smaller alphabet may mean only fewer candidates per step. Comparing accuracies over
    vocabularies of different sizes compares the difficulty of the task, not the quality of the
    method.
\end{enumerate}

A routing result is a claim about the route between units; a procedure result is a claim about
what gets done. With the alphabet declared beside each number there is nothing to correct; with
it omitted there is no way to correct.

One property of the source system constrains the rendering and belongs where the representation
is defined. A case records its actions in chronological order, but it can run in \emph{parallel}
across more than one organisational unit: once opened, it may be sent to several units at the
same time, each acting on it without the others having finished. A case's route is therefore not
a strictly linear queue of units but a set of stretches that can branch and reconverge. The
routing alphabet is a \emph{linearisation} of that route by the timestamp of each event, not a
reproduction of its concurrency structure. Two units acting simultaneously appear in the trace as
though one had followed the other, with the order decided by when each event was recorded. The
simplification is deliberate --- it keeps the trace in the sequence form the models of
Section~\ref{sec:typical-paths} require --- and its cost is stated in Section~\ref{sec:threats}.

\subsection{Typical-path models}\label{sec:typical-paths}

For each cluster we fit \textbf{six} generators, spread over three distinct ways of estimating
the next state. Two are \textbf{local counters}: a third-order Markov chain and a
directly-follows graph. Three are \textbf{neural sequence models}: GRU \citep{cho2014gru}, LSTM
\citep{hochreiter1997lstm} and Transformer \citep{vaswani2017attention}. The sixth is
\textbf{graph-convolutional} \citep{kipf2017gcn}, and generalises by a route none of the other
five uses. Clusters with fewer than 30 traces are not evaluated. Rolling a trained sequence model
forward to generate a whole trace, rather than scoring one-step predictions, is the arrangement
of \citet{camargo2019lstm}, and the realism of the generated sequence is what the path-support
measures of Section~\ref{sec:protocol} exist to gauge.

The three neural families are not chosen by architectural preference: they are the three lineages
that the next-activity literature consolidated --- the recurrent \citep{tax2017lstm}, the gated
recurrent, and the attention-based \citep{bukhsh2021processtransformer} --- and the reason to
compare all three rather than elect one is that under a common protocol no architecture dominates
across logs \citep{ramamaneiro2023dlppm}. Which one wins is therefore an empirical question
addressed to the corpus rather than a premise of the method. The three share one configuration,
fixed once and identical across every cluster and every evaluation cell with no per-cluster
tuning: state embeddings of 32 dimensions, a hidden state of 64, a context window of 8 states,
dropout 0.2, and 30 epochs of Adam at learning rate 0.01, weight decay $10^{-4}$, batches of
2{,}048 and a fixed seed. Fixing the configuration in advance is what stops the comparison
between families from measuring tuning effort instead of architecture.

The Markov generator is third-order with back-off: where a third-order context is unseen it falls
back to second and then first order, which is what keeps it usable on clusters near the 30-trace
floor. We use a directly-follows graph as one path learner --- first-order, restricted to the
training set's transitions --- while noting that this representation has documented pathologies
as a process model \citep{vanderaalst2019dfg}.

The graph-convolutional generator enters for a reason the other five do not cover. Markov and
directly-follows are strictly local estimators: the prediction for a state uses only the
transitions observed out of it, so a state seen three times in the partition has a three-sample
estimate and nothing corrects it. Over the training set's transition graph --- nodes are states,
edges are counted transitions, adjacency symmetrised, self-looped and symmetrically normalised
--- two rounds of propagation make a state's prediction incorporate its neighbours', so a rare
state borrows from structurally similar ones. It is the only family in the set that generalises
\emph{through the graph} rather than through the sequence, and it pays for that by conditioning on
the last state alone, with no context window: it gains neighbourhood and loses memory. It is
trained by cross-entropy over the observed pairs, with Adam, a hidden state of 64, 200 epochs,
learning rate 0.01, weight decay $10^{-4}$ and a fixed seed. Bit-for-bit reproducibility takes one
further precaution that is not obvious: linear-algebra parallelism is reduced to a single thread
during fitting, because sums in varying order diverge enough to swap the order of the top
predictions. Everything else --- prediction, next-state distribution, generation, beam, sampling
and constrained decoding --- is inherited from the counters, and the inheritance is deliberate:
the comparison measures the \emph{estimator}, not the decoder.

The directly-follows generator shares its estimator with one of the baselines. Under top-1
prediction it selects the most frequent observed successor, which is exactly the within-cluster
successor baseline of Section~\ref{sec:protocol}; the two are the same estimator and score
identically. There are therefore five distinct trained predictors plus a baseline that also
appears in the generator set, and the per-cluster maximum over ``six families'' is taken over a
set containing its own floor. Whether that overlap changes the selected-best figures is an
empirical question, answered in Section~\ref{sec:res-rq3}.

Generation is decoded under an automaton encoding the movement rules of the process system. What
the automaton guarantees should be stated at the strength the construction supports, because it
is a safety property rather than an acceptance condition. At every step the decoder admits only
states the automaton allows, so a produced path cannot receive a case that was never sent, nor
conclude a lane that was never opened. It does not guarantee that every send is matched before
the path terminates: termination is unconstrained, so a path may end with a transfer outstanding.
The constraint also binds only the six action types that participate in the movement rules ---
generation, remittance, receipt, reopening and the two conclusion types --- leaving the remaining
forty-two unconstrained. Section~\ref{sec:res-rq4} reports what the guarantee, so bounded,
actually buys.

\subsection{Evaluation protocol}\label{sec:protocol}

\subsubsection*{The unit of evaluation: the cell}

The elementary unit of evaluation is the \textbf{cell}: the crossing of a textual arm, a
parameterisation, a clustering algorithm and an alphabet. Inside a cell, each model family is
fitted \emph{separately in each cluster} and scored over the held-out traces of that same
cluster. A case enters a cell if it satisfies three chained conditions: it has a trace of at
least two states in that cell's alphabet; it has a label in that partition; and it belongs to a
cluster with at least thirty traces. The floor of thirty exists because a cluster split 70/30
needs enough evaluation cases for an accuracy to mean anything.

A reading trap follows from those conditions: the number of clusters a cell evaluates is
\textbf{not} the $k$ of the clustering that produced it. Clusters below the floor are excluded
from the evaluation without being excluded from the partition, and the first condition depends on
the alphabet --- so the same partition evaluates different populations in different alphabets.

\subsubsection*{The train/test split}

Cases are split 70/30 into training and held-out sets. The split is a deterministic hash of the
case identifier rather than a function of position in a list, so a case falls on the same side of
the split regardless of which subset it is drawn from --- within its own cluster, within the
global baseline, in any cell and in any alphabet. That property is what prevents leakage: the
global baseline trains on the union of the evaluated clusters' training sets, and because
membership does not depend on the list, no held-out trace can enter the pooled training set.
Without it the global baseline would be optimistic and the comparison worthless.

\subsubsection*{The measures, and the axis each one requires}

The primary measure is \textbf{next-step accuracy}, computed over \emph{decision steps} rather
than over traces: given the prefix of a held-out trace, what the next state is. A trace of $n$
states contributes $n-1$ steps, so long traces weigh more inside their cluster. Top-1 and top-3
are both reported. Beside it, generation quality is measured by an edit distance with
transpositions between the generated and the real suffix, conditioned on the true prefix of half
the trace --- the model receives what actually happened up to that point and completes the rest,
an arrangement in which lower is better.

High next-activity accuracy does not by itself demonstrate that process structure has been
captured \citep{peeperkorn2023rnnstructure}. We therefore report path-support measures alongside
accuracy: the fraction of units a real case visits that the generated path also visits, the
fraction of the generated path's units that real cases visit, the mean number of occurrences of
each generated transition per cluster member, and the count of transitions occurring in
\emph{no} sampled cluster member --- fabrication rather than variation. Complete generations are
also scored for \textbf{structural validity}, the degree of compliance with the movement rules the
decoder of Section~\ref{sec:typical-paths} enforces. Three further measures describe the partition
rather than the model and therefore take one value per cell, shared by all six families: the
fraction of held-out transitions already present in the training set's directly-follows graph, its
converse, and the harmonic mean of the two.

Two scope notes belong with those definitions. Transition support counts occurrences rather than
distinct members, so it may exceed 1 where a transition recurs within a case, and it is not a
proportion. Unseen transitions are counted against a sample of at most thirty held-out traces per
cluster rather than against the whole cluster, so the measure over-counts fabrication for clusters
larger than that. Appendix~\ref{app:glossary} defines every measure used.

One rule governs which of those measures a cell emits, and it follows directly from the alphabets:
\textbf{a measure is computed when, and only when, the axis it measures exists in that cell's
alphabet}. The action-by-unit alphabet carries both axes and emits everything. The procedure
alphabet carries only the action, so it emits structural validity and no unit measure. The routing
alphabet carries only the unit, so it emits the unit measures and \emph{no} structural validity.
Where the axis does not exist the key simply does not appear in the artefact. The alternative ---
computing the measure anyway --- returns not a wrong value but a value \emph{empty of referent}
that renders as legitimate: a movement constraint applied to states carrying no action never
binds, and would report perfect compliance by construction. Absence is honest; a maximum obtained
by vacuity is not.

\subsubsection*{Three baselines}

Three baselines, each answering a different question. The \textbf{within-cluster successor}
baseline predicts the most frequent next state inside the cluster, and asks whether the model
learned anything. The \textbf{global model} is a single third-order Markov chain trained on the
pooled training set of the whole corpus and evaluated on the same held-out traces; it asks
whether one model would have done as well. The \textbf{shuffled partition} permutes cluster
membership while preserving cluster count, size distribution and evaluated population; it asks
whether the grouping, rather than the shape of the partition, is doing the work.

Evaluating a clustering by the downstream model it enables is an established principle
\citep{deweerdt2013actitrac}. We extend it from process-model fitness to next-activity
prediction, apply it externally to partitions produced independently rather than inside a
clustering wrapper, and add the last two controls, which this literature does not routinely
report. The comparison against the global model is made with the model family fixed on both
sides --- per-cluster Markov chain against global Markov chain --- rather than by taking the best
of six families in each cluster against a single global family. Both forms are computed and
reported separately, so a reader can see what each choice of comparator produces.

\subsubsection*{Aggregation: two readings}

Each cluster produces its own accuracy; the cell aggregates those accuracies in two ways, and
both are reported. The \textbf{macro} reading is the simple mean over clusters, in which each
\emph{process type} counts the same however large. The \textbf{weighted} reading weights each
cluster by its number of traces, in which each \emph{case} counts the same. They are not two
versions of one quantity: they answer different administrative questions --- how well the method
serves the catalogue of types, and how well it serves the collection of cases.

The weighted reading carries an asymmetry. The weight is the cluster's number of traces, while
each cluster's accuracy is computed over \emph{decision steps}. A cluster of long traces
contributes many steps to its own accuracy but enters the aggregation only with its number of
cases. Where the length distribution varies sharply between clusters the two readings diverge,
and the divergence between them is the signal that it is happening.

\subsection{The estimand: gain against the control partition}\label{sec:estimand}

The quantity the study sets out to estimate is defined before any measurement, and the priority
counts as a commitment: it is against this that the results are read.

\subsubsection*{The partition dividend}

Fitting one model per cluster, rather than one model for the whole collection, changes what each
model faces. A cluster's effective alphabet is a subset of the global alphabet, and predicting the
next step within a smaller alphabet is easier; against that, each model now has a fraction of the
training data. Neither effect depends on the partition meaning anything: both follow only from how
many clusters there are and what sizes they have. Call the net effect of that change of shape the
\textbf{partition dividend}.

A gain measured against a global model therefore mixes two things: the dividend and whatever
cluster membership actually carries. Since the sign of the dividend is not known \emph{a priori}
--- it depends on which of the two effects predominates in the collection at hand --- a raw gain
cannot even be asserted to be a lower bound on the effect of interest.

\subsubsection*{The decomposition and the estimand}

Let $A(\Pi)$ be the aggregated accuracy of the per-cluster models under the real partition $\Pi$,
and $A(\tilde{\Pi})$ the same quantity under the shuffled control partition of
Section~\ref{sec:controls}, which preserves cluster count, size distribution and population. The
raw gain over the global model $B$ decomposes as

\[
\underbrace{A(\Pi) - B(\Pi)}_{\text{raw gain}} \;=\;
\underbrace{\left[A(\tilde{\Pi}) - B(\tilde{\Pi})\right]}_{\text{partition dividend}} \;+\;
\underbrace{\left[A(\Pi) - A(\tilde{\Pi})\right]}_{\text{partition signal}} \;+\;
\underbrace{\left[B(\tilde{\Pi}) - B(\Pi)\right]}_{\text{residual}} .
\]

\noindent \textbf{The estimand of this study is the second term.} For a model family $m$ and a
partition $\Pi$, with $A_w$ denoting the accuracy aggregated by each cluster's number of traces,

\[
\boxed{\;\Delta^{*}(m,\Pi) \;=\; A_w\!\left(m,\Pi\right) \;-\; A_w\!\left(m,\tilde{\Pi}\right)\;}
\]

\noindent and the adjusted variant, which subtracts the global model under both partitions,

\[
\Delta^{*}_{\mathrm{adj}}(m,\Pi) \;=\;
\left[A_w(m,\Pi) - B(\Pi)\right] \;-\; \left[A_w(m,\tilde{\Pi}) - B(\tilde{\Pi})\right] .
\]

\noindent The residual $B(\tilde{\Pi}) - B(\Pi)$ is the same global model scored over the same
held-out steps under both partitions, differing only through the weighting, and it therefore
serves as a free sanity check on the aggregation: an appreciable difference means the weighting is
doing work it should not. $\Delta^{*}$ is thus the difference between the real partition and a
shuffle of identical shape, and it is positive only if cluster membership carries information that
the shape of the partition does not.

\subsubsection*{\texorpdfstring{$\Delta^{*}$}{Delta-star} is an admission gate, not a ranking}

One property of $\Delta^{*}$ must be declared before any result, because it restricts the
quantity's legitimate use. The control preserves the shape of the partition but does not equalise
the \emph{marginal cost} of fragmenting the training data: a cluster formed at random has training
data with no internal structure, and its model degrades with size differently from a coherent
cluster's. $\Delta^{*}$ therefore grows with granularity --- a finer partition has more to gain
over its own shuffle, because specialising pays more when clusters are small --- and it is not
invariant to the number of clusters.

The legitimate use is as an \textbf{admission criterion}: a partition whose $\Delta^{*}$ is not
positive carries no signal beyond its shape and should not enter the study, whatever its internal
validity. To \emph{choose} among admitted partitions we use the explicit selection score of
Section~\ref{sec:setup}, which weighs downstream performance, coverage and quality. Comparisons at
matched cluster count remain valid under $\Delta^{*}$.

\subsubsection*{The two comparators answer different questions}

Neither the global model nor the shuffle is dispensable, because the questions are not the same.
The global model asks \emph{would a single model have done as well?} --- an engineering question,
whose answer guides whoever decides if a model per type is worth maintaining. The shuffle asks
\emph{does this grouping mean anything?} --- a validity question, whose answer guides whoever
decides if the catalogue is a real object. Both are computed and reported side by side, per
alphabet, and neither is presented as a substitute for the other.
